% Drawing 138/06-6 -- folder 138, batch 06, archive page 103 (PDF page 104).
% Transcription: un segment de $\times$ en $t = -3$ à $\bullet\,0$, puis une boucle issue de $0$ passant par $\times$ en $s = \frac32$ ; à l'intérieur de la boucle, le point $1$ (biffures) ; $\infty$ à l'extérieur.
% Kind: dessin d'enfant (segment and loop). Method 3 (encoded, rendered from the data). Status: ok.
% Notes: A segment from a cross at t = -3 to the bullet 0, then a loop from 0 through a cross at s = 3/2; inside the loop the point 1 (overwritten, as the transcription says), infinity outside. Encoded as dessin/1: three edges, the loop's two edges consecutive at 0; derived 1 black (degree 3), 2 white (degrees 1, 2), 2 faces (degrees 1 inside, 2 outside), genus 0. Notation of the page: a bullet over the value 0 is a black vertex, a cross over 1 a white vertex, a small circle over infinity the pole in a face; the number under each symbol is the multiplicity.
% Coordinates: 1 unit = 1 mm = 1/100 of the crop width (bbox in page fractions: [0.277, 0.794, 0.532, 0.833]).
% Redrawn by Claude (claude-opus-5-5), issue #73, 2026-11-07; not checked by a human; not compiled (no TeX in the session).
% Rights: reproduces part of the work -- behind RENDER_TRANSCRIPTS, never published by itself.
\documentclass[tikz,border=3pt]{standalone}
\usepackage{amsmath,amssymb}
\usetikzlibrary{patterns,arrows.meta}
\tikzset{pencil/.style={black!45}, struck/.style={black!35}, mark/.style={circle,draw,inner sep=1.3pt}, pt/.style={circle,fill,inner sep=0.9pt}, opt/.style={circle,draw,fill=white,inner sep=1.1pt}}
\begin{document}
\begin{tikzpicture}[x=1mm,y=1mm,line cap=round,line join=round,font=\scriptsize]
\draw (12.3,-12.3) -- (30.1,-12.3);
\draw (40.4,-11.6) ellipse[x radius=11, y radius=11];
\node[pt] at (30.1,-12.3) {};
\draw (10.8,-10.8) -- (13.8,-13.8);
\draw (10.8,-13.8) -- (13.8,-10.8);
\draw (50.5,-10.1) -- (53.6,-13.2);
\draw (50.5,-13.2) -- (53.6,-10.1);
\node[inner sep=1pt] at (12.3,-6.8) {\tiny $t=-3$};
\node[inner sep=1pt] at (25.3,-6.8) {\tiny $0$};
\node[inner sep=1pt] at (64.4,-7.5) {\tiny $s=\frac32$};
\node[inner sep=1pt] at (38.4,-7.5) {\tiny $1$};
\node[opt] at (38.4,-11.6) {};
\node[inner sep=1pt] at (91.8,-6.8) {\tiny $\infty$};
\node[opt] at (91.8,-11) {};
\end{tikzpicture}
\end{document}
